% ============================================================
%  Chapter 1 — Introduction
% ============================================================
\chapter{Introduction}
\label{chap:introduction}

\section{Motivation}
\label{sec:motivation}
% Begin with the high-level problem: modern codebases are complex, and AI coding assistants struggle with a lack of repository-level context.

\section{Problem Statement}
\label{sec:problem-statement}
% Narrow down to the specific problem your thesis addresses: creating an efficient and accurate method for retrieving context-relevant code snippets from an entire repository to feed into an LLM.

\section{Motivating Example}
\label{sec:motivating-example}
% Present a real, challenging example from your SWE-bench evaluation where a simple keyword search (like your BM25 baseline) fails, but your graph-based PPR approach succeeds.

\section{Proposed Solution and Contributions}
\label{sec:solution-contributions}
% Clearly state your solution at a high level. Provide a bulleted list of specific contributions.
Our contributions include:
\begin{itemize}
    \item The design and implementation of a novel graph-based context retrieval system using Personalized PageRank (PPR).
    \item A comprehensive evaluation on the SWE-bench Lite benchmark, demonstrating a significant improvement in retrieval recall over baseline methods.
    \item An analysis of the limitations of purely structural graph representations for code retrieval.
\end{itemize}

\section{Thesis Outline}
\label{sec:outline}
% Conclude with a brief overview of the rest of the paper's structure.
The remainder of this thesis is organized as follows:
\cref{chap:background} provides the necessary background and reviews related work.
\cref{chap:design} describes the system architecture and the retrieval algorithm.
\cref{chap:implementation} details the implementation of the \texttt{codegraph} system.
\cref{chap:evaluation} presents the experimental setup and results.
\cref{chap:discussion} discusses the findings and identifies limitations.
Finally, \cref{chap:conclusion} concludes the thesis and suggests future research directions.
